% Part: sets-functions-relations
% Chapter: functions
% Section: kinds

\documentclass[../../../include/open-logic-section]{subfiles}

\begin{document}

\olfileid[ar]{sfr}{fun}{kin}
\olsection{أنواع الدوال}

\begin{explain}
ولكثرة ورود بعض أصناف الدوال ننتفع بتصنيف يميزها.

فلنبدأ بما تبلغ قيمه كل عنصر من المجال المقابل. هذا الصنف
يسمى !!{surjective}، وصورته في \olref{fig:surjective}.

\begin{figure}
  \olasset{assets/diagrams/surjective.tikz}
  \caption{الدالة !!^a{surjective} تبلغ بقيمها كل !!{element} من المجال المقابل.}
  \ollabel{fig:surjective}
\end{figure}
\end{explain}

\begin{defn}[دالة !!^{surjective}]
الدالة $f \colon A \rightarrow B$ \emph{!!{surjective}} إذا وفقط إذا
اتحد $B$ ومدى~$f$؛ أي لكل $y \in B$ يوجد $x \in A$ واحد على الأقل
يحقق~$f(x) = y$. وصورته الرمزية:
\[
  (\forall y \in B)(\exists x \in A)f(x) = y.
\]
وتسمى مثل هذه الدالة !!a{surjection} من $A$ إلى $B$.
\end{defn}

\begin{explain}
فإثبات أن $f$ هي !!a{surjection} يكون بإظهار أن كل شيء في المجال
المقابل لـ$f$ يحصل قيمةً $f(x)$ لمدخل $x$ ما.

وكل دالة \emph{تستحث} !!a{surjection}. خذ $f \colon A \to B$،
واجعل $f' \colon A \to \ran{f}$ معرفة بـ$f'(x) = f(x)$.
فـ$\ran{f}$ \emph{بتعريفه} $\Setabs{f(x) \in B}{x \in A}$،
ومن ثم $f'$ هي !!a{surjection}.
\end{explain}

\begin{explain}
ولكل مدخل مخرج واحد في أي دالة. وقد تزيد الدالة على ذلك بأن
لا تجمع مدخلين مختلفين على مخرج واحد؛ فتسمى !!{injective}،
وتمثيلها في \olref{fig:injective}.
\begin{figure}
  \olasset{assets/diagrams/injective.tikz}
  \caption{الدالة !!^a{injective} تمنع اتفاق القيمة لوسيطتين مختلفتين.}
  \ollabel{fig:injective}
\end{figure}
\end{explain}

\begin{defn}[دالة !!^{injective}]
الدالة $f \colon A \rightarrow B$ \emph{!!{injective}} إذا وفقط إذا
كان لكل $y \in B$ عنصر $x \in A$ واحد على الأكثر يحقق~$f(x) = y$.
وتسمى حينئذ !!a{injection} من $A$ إلى~$B$.
\end{defn}

\begin{explain}
ولإثبات أن $f$ هي !!a{injection} خذ !!{element}s $x$ و$y$ من مجال
$f$، وبيّن أن $f(x)=f(y)$ لا يكون إلا مع $x=y$.
\end{explain}

\begin{ex}
خذ الثابتة $f\colon \Nat \to \Nat$ المحددة بـ$f(x) = 1$؛ فهذه
لا هي !!{injective} ولا هي !!{surjective}.

وأما الهوية $f\colon \Nat \to \Nat$، حيث $f(x) = x$،
فتجمع !!{injective} و!!{surjective}.

وأما اللاحق $f \colon \Nat \to \Nat$، حيث $f(x) = x+1$،
فهي !!{injective} دون !!{surjective}.

وللدالة $f \colon \Nat \to \Nat$ ذات التعريف:
\[
  f(x) =
  \begin{cases}
    \frac{x}{2} & \text{إذا كان $x$ زوجيًا} \\
    \frac{x+1}{2} & \text{إذا كان $x$ فرديًا.}
  \end{cases}
\]
صفة !!{surjective}، لا صفة !!{injective}.
\end{ex}

\begin{explain}
ويكثر احتياجنا إلى الدوال الجامعة لـ!!{injective} و!!{surjective}؛
واسمها !!{bijective}، وصورتها في \olref{fig:bijective}.
وتسمى !!^{bijection}s أيضًا \emph{مراسلات واحد إلى واحد}، إذ يقابل
كل عنصر من المجال المقابل فيها عنصر واحد من المجال لا غير.
\begin{figure}
  \olasset{assets/diagrams/bijective.tikz}
  \caption{الدالة !!^a{bijective} تجعل لكل عنصر من المجال المقابل مقابِلًا واحدًا في المجال.}
  \ollabel{fig:bijective}
\end{figure}
\end{explain}

\begin{defn}[!!^{bijection}]
الدالة $f \colon A \to B$ \emph{!!{bijective}} إذا وفقط إذا اجتمع
فيها !!{surjective} و!!{injective}. وحينئذ تسمى !!a{bijection}
من $A$ إلى~$B$، أو بين $A$ و~$B$.
\end{defn}

\end{document}
